\chapter{Conclusiones y Trabajo Futuro}
\label{cap:conclusiones}

Este capítulo resume las conclusiones principales del trabajo, discute las limitaciones del enfoque propuesto, y propone direcciones para investigación futura.

\section{Conclusiones}

\subsection{Sobre la Metodología Propuesta}

Este trabajo demostró la viabilidad de una metodología de testing basada en la reimplementación de subsistemas complejos de software. Los resultados principales son:

\subsubsection{Conclusión 1: La Reimplementación Permite Testing Controlado}

La reimplementación del módulo de detección de semáforos de Apollo en PyTorch permitió:

\begin{itemize}
    \item \textbf{Control total} sobre el código para instrumentación y debugging
    \item \textbf{Configuración simplificada} eliminando dependencias complejas (Docker, CyberRT)
    \item \textbf{Experimentación rápida} con modificaciones y casos de test específicos
    \item \textbf{Reproducibilidad} de tests en diferentes entornos
\end{itemize}

Este nivel de control habría sido imposible o extremadamente difícil trabajando directamente con el sistema Apollo completo.

\subsubsection{Conclusión 2: Reverse Engineering Sistemático es Factible}

El proceso de reverse engineering aplicado fue exitoso:

\begin{itemize}
    \item Se extrajo la especificación funcional completa del sistema
    \item Se documentaron todos los parámetros numéricos
    \item Se identificaron reglas de negocio no documentadas
    \item Se logró 100\% de equivalencia funcional validada
\end{itemize}

\textbf{Lección clave:} El análisis combinado (estático del código + dinámico con ejecuciones) fue esencial. Ninguno de los dos por separado habría sido suficiente.

\subsubsection{Conclusión 3: Oracle-Based Testing es Efectivo pero Insuficiente}

Los tests de validación (usando Apollo como oráculo) fueron exitosos:
\begin{itemize}
    \item 1150 tests de validación: 100\% passed
    \item Confirmaron equivalencia funcional
    \item Validaron cada etapa del pipeline independientemente
\end{itemize}

\textbf{Sin embargo:} Estos tests solo confirman que la reimplementación replica a Apollo. No identifican debilidades del diseño original.

\subsubsection{Conclusión 4: Tests Exploratorios Descubren Vulnerabilidades Reales}

Los tests exploratorios, diseñados con conocimiento del código interno, identificaron:

\begin{itemize}
    \item 3 vulnerabilidades de severidad media-alta
    \item 4 limitaciones funcionales significativas
    \item Múltiples casos límite no manejados explícitamente
\end{itemize}

\textbf{Hallazgos que no se habrían descubierto con black-box testing:}
\begin{enumerate}
    \item Asunción no validada de timestamps monotónicos
    \item Falta de validación de coordenadas de bboxes
    \item Degradación severa con imágenes sobreexpuestas
    \item Comportamiento indefinido con inputs adversariales específicos
\end{enumerate}

\subsubsection{Conclusión 5: El Enfoque es Generalizable}

La metodología propuesta no es específica a detección de semáforos:

\textbf{Requisitos para aplicabilidad:}
\begin{itemize}
    \item Subsistema con interfaces bien definidas
    \item Código fuente disponible (o posibilidad de reverse engineering)
    \item Posibilidad de generar/recolectar inputs de test
    \item Criterio para determinar correctitud (oráculo)
\end{itemize}

\textbf{Candidatos potenciales:}
\begin{itemize}
    \item Otros módulos de percepción (detección de peatones, carriles)
    \item Módulos de predicción de trayectorias
    \item Sistemas de planificación local
    \item Componentes de fusión de sensores
\end{itemize}

\subsection{Sobre el Sistema de Detección de Semáforos}

\subsubsection{Conclusión 6: El Diseño de Apollo es Generalmente Robusto}

El análisis reveló que Apollo está bien diseñado en aspectos fundamentales:

\begin{itemize}
    \item Parámetros numéricos bien calibrados
    \item Reglas de seguridad conservadoras (sesgo hacia evitar falsos negativos peligrosos)
    \item Arquitectura modular con separación clara de responsabilidades
    \item Manejo correcto de casos límite comunes
\end{itemize}

\subsubsection{Conclusión 7: Existen Limitaciones Inherentes al Enfoque}

Algunas limitaciones son inherentes al enfoque basado en proyecciones:

\begin{itemize}
    \item \textbf{Dependencia de HD-Maps:} Requiere mapas precisos actualizados
    \item \textbf{Solo detecta semáforos conocidos:} No descubre semáforos nuevos
    \item \textbf{Vulnerable a errores de localización:} Si la pose del vehículo es incorrecta, las proyecciones fallan
\end{itemize}

Estas limitaciones son trade-offs conscientes del diseño, no bugs.

\subsubsection{Conclusión 8: Hay Margen para Mejoras Específicas}

Los tests identificaron áreas concretas de mejora:

\begin{enumerate}
    \item \textbf{Validación de inputs:} Agregar checks de sanidad en coordenadas y timestamps
    \item \textbf{Robustez a iluminación:} Preprocesamiento adaptativo para manejar sobreexposición
    \item \textbf{Manejo de oclusión:} Considerar información temporal para inferir estado con oclusión severa
    \item \textbf{Reducción de falsos positivos:} Filtros adicionales para distinguir luces traseras de semáforos
\end{enumerate}

\section{Limitaciones del Trabajo}

\subsection{Limitaciones Metodológicas}

\subsubsection{Dependencia del Oráculo}

La validación asume que Apollo es correcto. Si Apollo tiene bugs, la reimplementación los replicará.

\textbf{Mitigación parcial:} Los tests exploratorios no dependen del oráculo y pueden descubrir bugs compartidos.

\subsubsection{Cobertura de Casos de Test}

Aunque se diseñaron 1150+ tests, no cubren exhaustivamente el espacio de entrada infinito.

\textbf{Casos potencialmente no cubiertos:}
\begin{itemize}
    \item Combinaciones específicas de condiciones adversariales
    \item Secuencias temporales complejas
    \item Casos extremos no anticipados
\end{itemize}

\subsubsection{Evaluación en Entorno Controlado}

Todos los tests se ejecutaron offline con datos pregrabados, no en un vehículo real.

\textbf{Factores del mundo real no evaluados:}
\begin{itemize}
    \item Vibración del vehículo
    \item Variación de sensores entre cámaras
    \item Latencia end-to-end en producción
    \item Interacción con otros módulos del sistema completo
\end{itemize}

\subsection{Limitaciones de Alcance}

\subsubsection{Un Solo Subsistema}

Este trabajo se enfocó únicamente en detección de semáforos. No se validó la metodología en otros subsistemas.

\subsubsection{Sistema Específico}

Se trabajó con Apollo específicamente. Otros sistemas de conducción autónoma pueden tener características diferentes que afecten la aplicabilidad del enfoque.

\subsection{Limitaciones Técnicas}

\subsubsection{Simplificaciones en la Reimplementación}

Algunas características de Apollo no se reimplementaron:
\begin{itemize}
    \item Proyección 3D→2D en tiempo real (se usaron projection boxes pre-calculadas)
    \item Soporte multi-cámara
    \item Integración con sistema de mapas HD completo
\end{itemize}

Estas simplificaciones no afectan la lógica core de detección/reconocimiento, pero limitan la evaluación del sistema completo.

\section{Trabajo Futuro}

\subsection{Extensiones Inmediatas}

\subsubsection{Implementación de Mejoras Identificadas}

Implementar las mejoras sugeridas por los tests:
\begin{enumerate}
    \item Agregar validación de inputs
    \item Implementar preprocesamiento adaptativo de contraste
    \item Mejorar manejo de oclusión con contexto temporal
    \item Reducir falsos positivos con luces traseras
\end{enumerate}

\textbf{Evaluación:} Medir mejora en métricas con nuevo dataset de test.

\subsubsection{Expansión de la Suite de Tests}

\begin{itemize}
    \item Aumentar cobertura de casos límite
    \item Agregar tests de regresión para las vulnerabilidades identificadas
    \item Diseñar tests de performance y eficiencia
    \item Crear tests de stress (e.g., 1000+ semáforos en una imagen)
\end{itemize}

\subsubsection{Evaluación con Más Datos}

\begin{itemize}
    \item Datasets públicos (LISA Traffic Light Dataset, Bosch TLD)
    \item Diferentes condiciones geográficas (otros países, tipos de semáforos)
    \item Condiciones climáticas adversas (lluvia, niebla, nieve)
    \item Video de larga duración para evaluación de tracking
\end{itemize}

\subsection{Investigación Metodológica}

\subsubsection{Generalización a Otros Subsistemas}

Aplicar la misma metodología a otros módulos de Apollo:

\begin{table}[H]
\centering
\begin{tabular}{|l|l|l|}
\hline
\textbf{Subsistema} & \textbf{Complejidad} & \textbf{Interés} \\
\hline
Detección de carriles & Media & Alto \\
Detección de peatones & Media-Alta & Muy Alto \\
Detección de obstáculos & Alta & Alto \\
Fusión de sensores & Muy Alta & Muy Alto \\
\hline
\end{tabular}
\caption{Candidatos para aplicar la metodología}
\end{table}

\textbf{Objetivo:} Validar que el enfoque funciona en diferentes contextos.

\subsubsection{Automatización del Proceso}

Desarrollar herramientas para automatizar partes del proceso:

\begin{enumerate}
    \item \textbf{Extracción automática de parámetros:} Parser de código C++ que identifica constantes numéricas
    \item \textbf{Generación automática de tests:} A partir de tipos de input/output
    \item \textbf{Fuzzing guiado:} Generación automática de casos adversariales
    \item \textbf{Análisis de cobertura:} Identificar qué partes del código no están cubiertas por tests
\end{enumerate}

\subsubsection{Formalización de Propiedades}

Desarrollar un lenguaje o framework para especificar formalmente propiedades del sistema:

\begin{codesnippet}
\begin{verbatim}
# Pseudo-código de especificación formal
@property
def no_red_to_green_transition(system):
    """Un semáforo nunca debe cambiar directamente de RED a GREEN"""
    for signal_id in system.tracked_signals:
        if signal_id.previous_state == RED:
            assert signal_id.current_state != GREEN or
                   signal_id.current_state == UNKNOWN
\end{verbatim}
\end{codesnippet}

\subsection{Mejoras al Sistema}

\subsubsection{Arquitecturas Alternativas}

Explorar arquitecturas de detección más modernas:
\begin{itemize}
    \item Detectores basados en Transformers (DETR, etc.)
    \item Arquitecturas más eficientes (EfficientDet, YOLO v8+)
    \item Modelos específicos para objetos pequeños
\end{itemize}

\textbf{Evaluación:} Comparar con la implementación de Apollo usando la misma suite de tests.

\subsubsection{Mejoras en Robustez}

\begin{enumerate}
    \item \textbf{Data augmentation durante entrenamiento:}
    \begin{itemize}
        \item Simulación de oclusión
        \item Variación de iluminación
        \item Ruido y blur
    \end{itemize}

    \item \textbf{Ensemble de modelos:}
    \begin{itemize}
        \item Múltiples detectores con votación
        \item Clasificadores entrenados con diferentes aumentaciones
    \end{itemize}

    \item \textbf{Uso de contexto temporal más sofisticado:}
    \begin{itemize}
        \item Redes recurrentes (LSTM) para tracking
        \item Filtro de Kalman para predicción de posición
        \item Modelos de lenguaje para secuencias de estados (e.g., Transformers)
    \end{itemize}
\end{enumerate}

\subsubsection{Integración de Incertidumbre}

Modificar el sistema para producir estimaciones de incertidumbre:

\begin{itemize}
    \item Clasificación con distribuciones (en lugar de clase única)
    \item Propagación de incertidumbre a través del pipeline
    \item Decisiones conservadoras cuando incertidumbre es alta
\end{itemize}

\textbf{Beneficio:} Permite toma de decisiones más segura en el módulo de planificación.

\subsection{Testing Avanzado}

\subsubsection{Adversarial Machine Learning}

Aplicar técnicas de adversarial ML específicamente:

\begin{enumerate}
    \item \textbf{FGSM (Fast Gradient Sign Method):} Generar perturbaciones mínimas que cambien clasificación
    \item \textbf{PGD (Projected Gradient Descent):} Ataques iterativos más fuertes
    \item \textbf{Physical adversarial patches:} Stickers que confunden al detector
\end{enumerate}

\textbf{Objetivo:} Identificar vulnerabilidades de seguridad (e.g., ataques adversariales en el mundo físico).

\subsubsection{Mutation Testing}

Aplicar mutation testing al código:
\begin{enumerate}
    \item Introducir mutaciones (cambios pequeños) en el código
    \item Ejecutar suite de tests
    \item Identificar mutaciones que sobreviven (tests no detectan el cambio)
    \item Mejorar tests para detectar esas mutaciones
\end{enumerate}

\textbf{Métrica:} Mutation score (% de mutaciones detectadas)

\subsubsection{Continuous Testing}

Establecer pipeline de testing continuo:
\begin{itemize}
    \item Ejecución automática de tests en cada commit
    \item Regresión visual (comparar imágenes de output)
    \item Benchmarking de performance
    \item Alertas ante degradación de métricas
\end{itemize}

\subsection{Aplicaciones Más Allá de Testing}

\subsubsection{Educación y Documentación}

La reimplementación puede servir como:
\begin{itemize}
    \item Material educativo para enseñar sistemas de percepción
    \item Documentación "ejecutable" del sistema Apollo
    \item Plataforma para experimentación académica
\end{itemize}

\subsubsection{Prototipado Rápido}

Usar la reimplementación como base para:
\begin{itemize}
    \item Probar nuevas ideas sin modificar Apollo completo
    \item Benchmarking de algoritmos alternativos
    \item Validación de conceptos antes de integración en sistema completo
\end{itemize}

\subsubsection{Certificación y Auditoría}

La metodología podría aplicarse para:
\begin{itemize}
    \item Auditoría independiente de sistemas de conducción autónoma
    \item Procesos de certificación regulatoria
    \item Validación de seguridad funcional (ISO 26262)
\end{itemize}

\section{Reflexiones Finales}

\subsection{Lecciones Aprendidas}

\subsubsection{Sobre Reverse Engineering}

\begin{itemize}
    \item La documentación del código a menudo es incompleta o desactualizada
    \item Muchas decisiones de diseño no están explícitas en el código
    \item El análisis dinámico (ejecución + observación) es tan importante como el análisis estático
    \item Validar cada asunción es crítico
\end{itemize}

\subsubsection{Sobre Testing de Sistemas ML}

\begin{itemize}
    \item Tests tradicionales de software son necesarios pero no suficientes
    \item Property-based testing es particularmente útil para sistemas ML
    \item Los casos límite más interesantes surgen del análisis del código
    \item No existe un "conjunto completo" de tests - siempre hay más casos por descubrir
\end{itemize}

\subsubsection{Sobre Sistemas de Conducción Autónoma}

\begin{itemize}
    \item La complejidad de estos sistemas justifica enfoques creativos de testing
    \item El acoplamiento entre componentes dificulta el testing aislado
    \item Las asunciones implícitas son una fuente común de bugs
    \item La seguridad requiere múltiples capas de validación
\end{itemize}

\subsection{Impacto Potencial}

Si esta metodología se adopta más ampliamente, podría:

\begin{enumerate}
    \item \textbf{Mejorar la calidad} de sistemas de conducción autónoma mediante testing más exhaustivo

    \item \textbf{Acelerar el desarrollo} permitiendo prototipado e iteración rápida

    \item \textbf{Facilitar la investigación} proporcionando implementaciones de referencia accesibles

    \item \textbf{Aumentar la transparencia} mediante documentación ejecutable de sistemas complejos

    \item \textbf{Contribuir a la seguridad} identificando vulnerabilidades antes del despliegue
\end{enumerate}

\subsection{Conclusión Final}

Este trabajo demostró que la reimplementación dirigida de subsistemas complejos es una estrategia viable y valiosa para el testing de sistemas de conducción autónoma. La combinación de reverse engineering sistemático, oracle-based testing, y tests exploratorios guiados por conocimiento del código interno permitió:

\begin{itemize}
    \item Validar equivalencia funcional con el sistema de referencia
    \item Identificar vulnerabilidades y limitaciones específicas
    \item Proponer mejoras concretas basadas en evidencia
\end{itemize}

El enfoque propuesto es generalizable a otros subsistemas y sistemas, representando una contribución metodológica al campo de testing de sistemas basados en machine learning para aplicaciones críticas de seguridad.

\vspace{1cm}

\textit{``Testing shows the presence, not the absence of bugs.''} - Edsger W. Dijkstra

\vspace{0.5cm}

Este trabajo aplicó esa sabiduría diseñando tests específicamente para descubrir bugs. Los hallazgos confirman que incluso sistemas bien diseñados como Apollo tienen margen para mejora - y que el testing creativo e informado es esencial para encontrar esas oportunidades.